\section{Introduction}
Standard-cell library characterization is a fundamental component of modern VLSI design flows, since the quality of cell timing models directly affects the accuracy and efficiency of static timing analysis (STA) \cite{bhasker2009sta}. In industrial practice, timing libraries are typically generated through extensive transistor-level simulation under a large number of input slew and output load conditions \cite{garyfallou2021gate,charafeddine2017curve}. As technology scales toward advanced nodes, however, library characterization becomes increasingly expensive due to stronger parasitic effects, more complex device behaviors, and the significantly higher cost of SPICE-level simulation \cite{amelifard2008csm,kaur2014ecsm}. Meanwhile, a practical standard-cell library usually contains hundreds to thousands of cells with multiple logic functions and drive strengths, making full characterization extremely time-consuming and creating a major bottleneck for library development and design technology co-optimization \cite{ma2024gnn,liu2025gtncell}.

To alleviate this cost, recent studies have explored machine learning based timing modeling for standard-cell characterization. Prior works have applied regression models, multilayer perceptrons, and graph-based neural networks to predict delay, power, or reliability-related timing behavior, demonstrating that learning-based approaches can effectively approximate circuit behavior while significantly reducing simulation overhead \cite{klemme2021reliability,garyfallou2022ml,ebrahimipour2020aadam}. Among them, graph neural networks are particularly attractive because they can naturally encode circuit topology and device-level interactions. Recent advances have further shown that graph attention and graph transformer based models can improve characterization accuracy by capturing richer structural dependencies inside standard cells \cite{ma2024gnn,ye2024aging,liu2025gtncell,velivckovic2018gat}. Moreover, parasitic-aware heterogeneous graph modeling has been shown to be promising for timing library characterization in advanced technologies \cite{cheng2024hgat}. Nevertheless, most existing methods are developed within a single technology node, and therefore rely on the availability of sufficient labeled data in the target process.

This assumption is often unrealistic for advanced technology nodes. Compared with mature planar technologies, advanced nodes usually incur much higher characterization cost per sample, while the amount of labeled target-node data is much more limited \cite{amuru2024pvt,zhang2024disentangle}. As a result, directly training a high-capacity predictor on the target node is often impractical. A natural solution is to transfer knowledge from mature nodes with abundant characterization data to advanced nodes with scarce data. Transfer learning has been widely studied as an effective paradigm for low-resource learning and domain adaptation \cite{weiss2016survey,zhuang2021survey,dai2009eigentransfer,ganin2016dann}. In circuit and timing prediction tasks, several recent studies have shown that cross-node transfer can reduce data requirements and improve modeling efficiency \cite{shrestha2024arrival,amuru2024pvt,zhang2025prerouting}. However, directly reusing a source-domain model or naively mixing source-node and target-node samples can suffer from severe domain shift \cite{zhang2024disentangle}. Across process nodes, the logical topology of the same standard cell is often largely preserved, whereas physical properties such as parasitics, threshold voltages, and process-dependent electrical behaviors may differ substantially. If such shared and domain-specific factors are not properly separated, negative transfer may occur and degrade prediction accuracy \cite{tang2024semisupervised,tang2025knowledge,zhang2024disentangle}.

Motivated by these observations, this paper proposes a cross-technology-node timing prediction framework for standard-cell libraries based on heterogeneous graph learning and transfer learning. Instead of representing a standard cell only at the gate level, we construct a heterogeneous circuit graph that explicitly models transistors, resistors, capacitors, ports, and their connectivity, so that both structural and parasitic information can be preserved. On top of this representation, we develop an H-GAT based encoder to capture heterogeneous interactions inside each cell through structure-guided attention and physics-aware message aggregation. Our design is inspired by recent heterogeneous graph attention mechanisms and parasitic-aware timing graph modeling \cite{yang2021hgat,cheng2024hgat}. Furthermore, to improve transferability across process nodes, we disentangle the learned representation into technology-invariant structural features and technology-dependent physical features. The former captures logical topology and reusable design knowledge shared across nodes, while the latter models process-related variations that require domain adaptation. Based on this decomposition, we introduce a transfer strategy that preserves shared structural knowledge while aligning physical representations across domains under low-resource target-node settings.

The proposed framework targets the practical scenario where abundant source-domain data from mature nodes and only a small amount of target-domain data from advanced nodes are available. In this work, we focus on cross-node timing prediction for standard-cell libraries using open PDKs such as Nangate45 and ASAP7 as representative source and target domains. Our goal is to provide an accurate and low-cost alternative to exhaustive HSPICE-based characterization, thereby improving the scalability of timing library development for advanced processes.

The main contributions of this work are summarized as follows:
\begin{itemize}

    \item We formulate cross-technology-node standard-cell timing prediction as a transfer learning problem under target-domain data scarcity, and identify the key challenge of separating shared circuit structure from process-dependent physical variation.
    \item We propose a heterogeneous graph representation for standard cells that explicitly incorporates circuit topology and parasitic components, enabling more faithful modeling of advanced-node electrical behavior.
    \item We design an H-GAT based timing encoder with structure-guided attention and physics-aware message aggregation, and further develop a disentanglement-and-alignment transfer strategy for cross-node adaptation.
    \item We establish a practical evaluation setting for cross-technology-node timing prediction with open PDKs, and compare the proposed method against target-only training, mixed-domain training, and pre-training plus fine-tuning baselines.
\end{itemize}
